\documentclass[11pt]{article}
\usepackage[margin=1in]{geometry}
\usepackage{amsmath,amssymb,amsthm}
\usepackage{algorithm}
\usepackage{algpseudocode}
\newtheorem{proposition}{Proposition}
\newtheorem{lemma}{Lemma}
\newtheorem{remark}{Remark}
\newtheorem{definition}{Definition}
\newcommand{\F}{\mathcal{F}}
\newcommand{\E}{\mathbb{E}}
\newcommand{\Prob}{\mathbb{P}}
\newcommand{\KL}{\mathrm{KL}}
\newcommand{\M}{\mathrm{M}}
\title{FX-Mix: fact-free exclusion distillation on native-sampler states\\for masked diffusion language models (LLaDA)\\[4pt]\large Principal definitions, algorithm and proofs (llada\_method\_design\_v1, SPEC v1.3.1)}
\author{XXResearch llada\_method\_design\_v1 sprint}
\date{2026-10-20}
\begin{document}
\maketitle

\paragraph{Status of the mathematics.} Lemma~\ref{lem:T} is the standard information (I-)projection onto a
probability constraint (Csisz\'ar), the target already used by TNT/TNFF and, as a limit, by UNDIAL. Proposition~\ref{prop:hazard}
is an elementary first-passage decomposition plus a union bound, instantiated for the pinned LLaDA sampler.
Proposition~\ref{prop:escape} is an elementary identifiability statement in the unrestricted (zero parameter sharing)
model class. None of these is claimed as the method's novelty, none connects the optimised loss to leakage of a
network with shared parameters, and none implies global forgetting or utility preservation.

\section{Setting}
\begin{definition}[Mask predictor and states]
Vocabulary $V$, mask symbol $m\notin V$. A prompt $c$ is never masked. The answer region has fixed length $L$
(\texttt{gen\_length}, EOS padding is ordinary vocabulary). A \emph{state} is $s=(c,z)$ with $z\in(V\cup\{m\})^L$ and
masked set $\M(s)=\{i: z_i=m\}$. A mask predictor maps $(s,i)$, $i\in\M(s)$, to logits $\zeta(s,i)\in\mathbb{R}^{V}$ and
$p(\cdot\mid s,i)=\mathrm{softmax}\,\zeta(s,i)$. $p_\theta$ denotes the student, $p_0=p_{\theta_0}$ the frozen reference.
\end{definition}
LLaDA's SFT objective draws $t\sim U(0,1)$, masks each answer token independently with probability $(1-\epsilon_t)t+\epsilon_t$
($\epsilon_t=10^{-3}$) and minimises $\sum_{i\in\M}\!-\log p_\theta(y_i\mid s,i)/p_{\text{mask}}$ divided by the SFT answer length (incl.\ EOS padding); it
therefore trains only on \emph{teacher-forced} states whose visible answer tokens agree with the reference answer.

\begin{definition}[Native sampler class $\Pi_{\le 1}$; pinned \texttt{generate} semantics]\label{def:sampler}
A configuration $\pi=(N,b,\text{rule},T)$ with $b\mid L$, $(L/b)\mid N$, rule $\in\{\mathrm{lc},\mathrm{rnd}\}$,
$T\in[0,1]$, \texttt{cfg\_scale}$=0$, no logit edits before the draw (\texttt{logits\_eos\_inf}=False, no end-think boost).
Start from the all-mask answer. At step $n$ (in block $\beta(n)$) let $s_n$ be the current state and
$E_n(s_n)=\M(s_n)\cap\beta(n)$ the eligible positions; draw $x^0_i\sim\pi_T(\cdot\mid s_n,i)=\mathrm{softmax}(\zeta(s_n,i)/T)$
independently over $i$ (Gumbel-max; $T=0$: arg-max with the implementation's fixed tie-break); select $S_n\subseteq E_n$
with $|S_n|=k_n$ (deterministic uniform schedule): the $k_n$ largest $p(x^0_i\mid s_n,i)$ (lc) or a uniformly random
$k_n$-subset independent of the draws (rnd); commit $x^0_i$ for $i\in S_n$; never re-mask.
\end{definition}

\begin{definition}[Forget request, leakage, fact-free states]
A forget item is $(c,y^\star,\F)$ with $\F\subset V$ nonempty (fact-token set, defined per fact) and no token of $\F$ in $c$.
$\tau=\min\{n: \exists i\in S_n,\ x^0_i\in\F\}$ ($\tau=\infty$ if none); the leakage event is $\Lambda_\F=\{\tau<\infty\}$,
counting all $L$ positions (including those after an EOS). A state is \emph{fact-free} if no visible answer token lies in $\F$.
$\rho(s,i)=p(\F\mid s,i)/\max_{v}p(v\mid s,i)$.
\end{definition}

\section{Exclusion target}
\begin{lemma}[Exclusion target; standard I-projection]\label{lem:T}
Let $p_0=\mathrm{softmax}(\zeta_0)$ with $p_0(\F)<1$, $\varepsilon\in[0,1)$ and
$r_\varepsilon=\arg\min_{q:\,q(\F)\le\varepsilon}\KL(q\,\|\,p_0)$. Write $p_0|_A$ for $p_0$ conditioned on $A\subseteq V$.
\begin{enumerate}
\item[(i)] $r_\varepsilon=p_0$ if $p_0(\F)\le\varepsilon$; otherwise $r_\varepsilon=\varepsilon\,p_0|_\F+(1-\varepsilon)\,p_0|_{\F^c}$,
for $\varepsilon>0$ equivalently $\mathrm{softmax}(\zeta_0-\gamma\mathbf{1}_\F)$ with $\gamma=\log\frac{p_0(\F)(1-\varepsilon)}{\varepsilon(1-p_0(\F))}$
(for computation use fp32 logits and the log-space form $\log r_\varepsilon(v)=\log(1-\varepsilon)+\zeta_{0,v}-\mathrm{LSE}_{\F^c}\zeta_0$ for $v\notin\F$,
$\log\varepsilon+\zeta_{0,v}-\mathrm{LSE}_{\F}\zeta_0$ for $v\in\F$; $\varepsilon=0$ gives $r_0=p_0|_{\F^c}$).
\item[(ii)] $\KL(r_\varepsilon\|p_0)=\mathrm{kl}(\varepsilon\,\|\,p_0(\F))$ if $p_0(\F)>\varepsilon$ (else $0$), where
$\mathrm{kl}$ is the binary KL; for $\varepsilon=0$ this equals $-\log(1-p_0(\F))$.
\item[(iii)] $\nabla_\zeta \KL(r_\varepsilon\|\mathrm{softmax}\,\zeta)\big|_{\zeta_0}=(p_0(\F)-\varepsilon)_+\,(p_0|_\F-p_0|_{\F^c})$,
which vanishes iff $p_0(\F)\le\varepsilon$.
\item[(iv)] For every distribution $p$ on $V$: $\KL(r_\varepsilon\|p)\le\kappa\Rightarrow \mathrm{kl}(r_\varepsilon(\F)\|p(\F))\le\kappa$; hence
$p(\F)\le x_\varepsilon(\kappa):=\max\{x:\mathrm{kl}(\varepsilon\|x)\le\kappa\}$ when $r_\varepsilon(\F)=\varepsilon$, and $x_0(\kappa)=1-e^{-\kappa}$.
\end{enumerate}
\end{lemma}
\begin{proof}
(i,ii) Chain rule over the partition $\{\F,\F^c\}$: $\KL(q\|p_0)=\mathrm{kl}(q(\F)\|p_0(\F))+q(\F)\KL(q|_\F\|p_0|_\F)+q(\F^c)\KL(q|_{\F^c}\|p_0|_{\F^c})$.
The last two terms are minimised (to $0$) by $q|_A=p_0|_A$, and $x\mapsto\mathrm{kl}(x\|p_0(\F))$ is convex with minimum at
$p_0(\F)$, so on $[0,\varepsilon]$ it is minimised at $\min(\varepsilon,p_0(\F))$. The softmax form follows by direct substitution.
(iii) $\nabla_\zeta\KL(r\|\mathrm{softmax}\,\zeta)=\mathrm{softmax}(\zeta)-r$; at $\zeta_0$, for $v\in\F$:
$p_{0,v}-\varepsilon p_{0,v}/p_0(\F)=(p_0(\F)-\varepsilon)(p_0|_\F)_v$; for $v\notin\F$:
$p_{0,v}-(1-\varepsilon)p_{0,v}/(1-p_0(\F))=-(p_0(\F)-\varepsilon)(p_0|_{\F^c})_v$.
(iv) Data processing under the coarse-graining $V\to\{\F,\F^c\}$; $\mathrm{kl}(\varepsilon\|\cdot)$ is increasing on $[\varepsilon,1)$.
\end{proof}

\section{Algorithm}
\begin{algorithm}[h]
\caption{FX-Mix (one optimisation step; SPEC v1.3.1 \S2)}
\begin{algorithmic}[1]
\Require student $\theta$ (LoRA, $B{=}0$ init), frozen $\theta_0$ (adapter off), forget items $\{(c,y^\star,\F)\}$, retain set,
pool $\Pi_{\text{train}}=\{\pi_k\}$ with weights $w_k$, $\alpha,\varepsilon,J,R,\lambda_f,\lambda_r$
\If{step $\equiv 0 \pmod R$} \Comment{rollout pool, no gradient}
  \For{each forget item} take $\pi_k$ from the pool (implementation: stratified round-robin over forget prompts, each config for $1/|\Pi_{\text{train}}|$ of them); run \texttt{generate}$(\theta,\pi_k)$; record $s_n$ before each step;
  stop at $\tau$; store $J$ pre-$\tau$ states sampled uniformly, plus $s_\tau$ if $\tau<\infty$
  \EndFor
\EndIf
\State $\mathcal{S}\gets$ fraction $\alpha$: TF$_{\text{ff}}$ states ($t\sim U(0,1)$; mask non-fact positions w.p.\ $t$; mask all $j$ with $y^\star_j\in\F$);
fraction $1-\alpha$: states from the pool
\State $\ell_f\gets \frac{1}{|\mathcal S|}\sum_{s\in\mathcal S}\frac{1}{|\M(s)|}\sum_{i\in\M(s)}\KL\big(r_\varepsilon(\cdot\mid s,i)\,\|\,p_\theta(\cdot\mid s,i)\big)$,
$r_\varepsilon$ from Lemma~\ref{lem:T} applied to $p_0(\cdot\mid s,i)$ (one reference pass, log-space)
\State $\ell\gets\lambda_f\ell_f+\lambda_r\,\mathcal{L}_{\text{SFT}}(\theta;\text{retain batch})$; AdamW (weight decay 0) step with clipping $1.0$
\end{algorithmic}
\end{algorithm}
\noindent The per-rollout $J$-sampling is a design measure, not an unbiased sample of the visitation measure. The average over positions in $\ell_f$ is a design choice; it is \emph{not} the integrand of
Proposition~\ref{prop:hazard}, and $\ell_f$ is not a leakage bound.

\section{Leakage under native samplers}
\begin{proposition}[First-commit decomposition and sampler-uniform hazard bound]\label{prop:hazard}
Fix $c,\F$ and a deterministic model. For every $\pi\in\Pi_{\le1}$:
\begin{enumerate}
\item[(a)] $\Prob(\Lambda_\F)=\sum_{n}\E\big[\mathbf{1}\{\tau\ge n\}\,h_n(s_n)\big]$ with $h_n(s)=\Prob(\exists i\in S_n: x^0_i\in\F\mid s_n=s)$,
and $s_n$ is fact-free on $\{\tau\ge n\}$.
\item[(b)] $h_n(s)\le\sum_{i\in E_n(s)}\pi_T(\F\mid s,i)$ ($h_n=0$ when $k_n=0$, in which case $E_n=\emptyset$); for rule $=\mathrm{rnd}$, $h_n(s)\le\frac{k_n}{|E_n(s)|}\sum_{i\in E_n(s)}\pi_T(\F\mid s,i)$.
\item[(c)] $\pi_T(\F\mid s,i)\le\min\{1,\rho(s,i)\}$ for all $T\in[0,1]$, and $\pi_1(\F\mid s,i)=p(\F\mid s,i)$.
\item[(d)] Hence $\Prob(\Lambda_\F)\le B_\pi:=\E\big[\sum_{n\le\tau\wedge(N-1)}\sum_{i\in E_n(s_n)}\min\{1,\rho(s_n,i)\}\big]$
(with the factor $k_n/|E_n|$ inside for rnd).
\item[(e)] If $T=0$ and $\rho(s,i)<1$ for every fact-free state $s$ reachable under $\pi$ and every $i\in E_n(s)$, then
$\Prob(\Lambda_\F)=0$ for any tie-break rule.
\end{enumerate}
\end{proposition}
\begin{proof}
(a) The events $\{\tau=n\}$ are disjoint and their union is $\Lambda_\F$. Given the history, step $n$ uses fresh
Gumbel/uniform noise and a deterministic model, so $\Prob(\tau=n\mid\text{history})=\mathbf{1}\{\tau\ge n\}h_n(s_n)$;
on $\{\tau\ge n\}$ no $\F$ token has been committed, so $s_n$ is fact-free.
(b) $\{\exists i\in S_n:x^0_i\in\F\}\subseteq\bigcup_{i\in E_n}\{x^0_i\in\F\}$, which holds even though lc selection depends on the draws;
union bound. For rnd, $S_n$ is a uniform $k_n$-subset independent of $x^0$, so $\Prob(i\in S_n)=k_n/|E_n|$.
(c) $T=0$: an $\F$ arg-max requires $p(\F)\ge\max_v p(v)$, i.e.\ $\rho\ge1$. $T\in(0,1]$, $\beta=1/T\ge1$:
$\sum_{f\in\F}p_f^\beta/\sum_v p_v^\beta\le\sum_f p_f^\beta/p_{\max}^\beta\le\sum_f p_f/p_{\max}=\rho$.
(d) Combine (a)--(c). (e) $\rho<1$ implies $p_f<p_{\max}$ strictly for every $f\in\F$, so no $\F$ token is an arg-max.
\end{proof}
\noindent\textbf{Use.} $\hat B_\pi$, the Monte-Carlo average over fresh rollouts of the \emph{final} model of the eligible-sum
bracket $\sum_{n\le\tau\wedge(N-1)}\sum_{i\in E_n(s_n)}\pi_T(\F\mid s_n,i)$ of (b) with the exact tempered probability of the evaluated
configuration ($\times k_n/|E_n|$ for rnd; $T=0$: $\mathbf{1}\{\arg\max\in\F\}$), in
deployment precision, is an unbiased estimate (report its Monte-Carlo standard error) of an upper bound on that model's leakage under $\pi$ (post-hoc
certificate). \textbf{Boundary cases} (all verified numerically in \texttt{agents/E\_scripts}): $T>1$ breaks (c)
($T=2$: $0.25>0.11$); \texttt{cfg\_scale}$>0$ can raise $p(\F)$ from $0.05$ to $0.99$ (scale 2, unconditional fact mass $0.001$); \texttt{logits\_eos\_inf} violates the sampler definition and breaks (c,e) ($\rho=0.5$, yet greedy commits $f$);
$\rho=1$ ties are resolved to the lowest id by \texttt{torch.argmax}; bf16 confidence ties make greedy paths device-dependent.
The bound is tight when all eligible positions are committed at once, $\approx2\times$ loose for rnd without the factor,
and $\approx9\times$ (eligible-sum form) to $16\times$ ($\rho$ form) loose for lc at $T=1$ after training (B's tabular study).

\section{Why teacher-forced states do not identify config-robust forgetting}
Let $\mathcal{S}_{\mathrm{ff}}$ be the fact-free states of a forget item, $r=r_0$ the exclusion target (Lemma~\ref{lem:T}, $\varepsilon=0$)
of the reference at each $(s,i)$, and $\bar r$ the model that equals $r$ at every fact-free state (``decode-time ban'' of the reference).
For $S\subseteq\mathcal{S}_{\mathrm{ff}}$ let $\mathcal{P}_S$ be the \emph{unrestricted} class of models whose conditionals equal $r$ at
every $s\in S$ and are arbitrary elsewhere (no parameter sharing), and define the escape time $\sigma_S=\min\{n\le N-1:\ s_n\notin S,\ \M(s_n)\neq\emptyset\}$
and the escape probability $e_\pi(S)=\Prob^{\bar r}_\pi(\sigma_S<\infty)$. (States change only through commits and every
masked position is eventually committed, so if $\M(s_{\sigma_S})\neq\emptyset$ the next commit is made from $s_{\sigma_S}$.)
\begin{proposition}[Escape characterisation; zero-sharing class]\label{prop:escape}
Fix $\pi\in\Pi_{\le1}$ and $S\subseteq\mathcal{S}_{\mathrm{ff}}$.
\begin{enumerate}
\item[(a)] For every $p\in\mathcal{P}_S$, $\Prob^p_\pi(\Lambda_\F)\le e_\pi(S)$, and $\sup_{p\in\mathcal{P}_S}\Prob^p_\pi(\Lambda_\F)=e_\pi(S)$;
$\inf_{p\in\mathcal{P}_S}\Prob^p_\pi(\Lambda_\F)=0$ (attained by $\bar r$).
\item[(b)] If $S$ contains every state visited with positive probability by $\pi$ under $\bar r$, then $e_\pi(S)=0$ and every
$p\in\mathcal{P}_S$ has $\Prob^p_\pi(\Lambda_\F)=0$.
\item[(c)] Let $S_{\mathrm{TF}}$ be the teacher-forced fact-free states (every visible answer token equals $y^\star$, every
position $j$ with $y^\star_j\in\F$ masked), and let $m\ge1$ be the number of such fact positions. Under $\bar r$ every commit at a fact
position differs from $y^\star_j$, so the trajectory escapes $S_{\mathrm{TF}}$ no later than the state following the first commit
at a fact position, unless no masked position remains (it escapes earlier whenever a non-fact position receives a token $\neq y^\star_j$). In particular, for one-token-per-step rnd decoding ($N=L$, $b=L$),
$e_\pi(S_{\mathrm{TF}})\ge 1-\frac{1}{L}$ if $m=1$ and $e_\pi(S_{\mathrm{TF}})=1$ if $m\ge2$, irrespective of $T$ and of $p_0$.
\item[(d)] For $\varepsilon>0$ and $\mathcal{P}^\varepsilon_S$ (conditionals $=r_\varepsilon$ on $S$), the same argument gives
$\sup_{p\in\mathcal{P}^\varepsilon_S}\Prob^p_\pi(\Lambda_\F)=\Prob^{\bar r_\varepsilon}_\pi(\tau<\sigma_S\ \text{or}\ \sigma_S<\infty)$, since leakage before
escape is then possible; the method's default $\varepsilon>0$ is covered only by this item.
\end{enumerate}
\end{proposition}
\begin{proof}
(a) For $n<\sigma_S$, either $s_n\in S$, where $p=r$, or $\M(s_n)=\emptyset$ and nothing is committed; at states in $S$, $r(\F)=0$, so no $\F$ token can be committed before $\sigma_S$,
and the law of the trajectory up to $\sigma_S$ is the same under $p$ and $\bar r$ (the draws and the lc confidences $p(x^0_i\mid s_n,i)$, $i\in E_n$,
depend only on the conditionals at the current state; the noise, the schedule $k_n$ and the tie-break are model-independent). Hence $\Lambda_\F\subseteq\{\sigma_S<\infty\}$ and the probability of the
latter is $e_\pi(S)$ for every $p\in\mathcal{P}_S$. For the supremum, take $p_\delta$ equal to $r$ on $S$ and with logits
$\zeta(s,i)=\delta^{-1}\mathbf{1}_{f}$ for a fixed $f\in\F$ at every fact-free $s\notin S$; the state $s_{\sigma_S}$ is fact-free
(no $\F$ token was committed before) and has masked positions, so the next step with $k_n\ge1$ commits from $s_{\sigma_S}$, where
every eligible draw equals $f$ with probability $\to1$ as $\delta\to0$ (for $T=0$, exactly); hence
$\Prob^{p_\delta}(\Lambda_\F)\to e_\pi(S)$. The infimum: $\bar r$ never commits an $\F$ token.
(b) Immediate from (a).
(c) Under $\bar r$, $r(y^\star_j\mid s,j)=0$ at fact positions because $y^\star_j\in\F$, so a committed fact position holds a
token $\ne y^\star_j$ and the next state (if masked positions remain) is outside $S_{\mathrm{TF}}$. With one token per step in
uniformly random order (rnd, $b=L$), the positions are committed in a uniformly random permutation independent of the draws;
no escape requires every fact position to be committed at the final commit step (and every earlier commit to equal $y^\star$), which has probability at most $1/L$ for $m=1$ and $0$ for
$m\ge2$.
\end{proof}
\begin{remark}[Scope of Proposition~\ref{prop:escape}]
It is a statement about the unrestricted class: it shows that exact agreement with the target on teacher-forced states does
not \emph{identify} leakage under random-order decoding, while agreement on the target sampler's own visitation does. With
parameter sharing nothing is proved either way; whether a network trained on $S_{\mathrm{TF}}$ generalises to the escaped
states is the empirical question tested by the SC-TF vs FX-Mix comparison. In the locked CPU validation (RESULTS.md \S3) the answer
was negative for the toy network: SC-TF matched FX-Mix on held-out leakage although $86\%$ of SC-TF's residual first-commit
decision states lay off the teacher-forced manifold, i.e.\ parameter sharing transferred the suppression. Part (b) is self-referential: for $p\in\mathcal{P}_S$ the laws agree up to $\sigma_S$, so $S$ contains every state $p$ visits under
$\pi$ iff it contains every state $\bar r$ visits. A student that matches the target on its \emph{own} visitation therefore does not leak
under $\pi$ (DAgger-type fixed point). FX-Mix approximates this with a finite, stale pool and $\varepsilon>0$; a finite pool gives no
guarantee for states of small visitation probability.
\end{remark}

\section{Remarks (not propositions)}
\begin{remark}[Consistent joint vs.\ per-position target]
For a consistent joint $p_0(y\mid c)$, $\arg\min_{q:q(\Lambda)=0}\KL(q\|p_0)=p_0(\cdot\mid\Lambda^c)$ with value $-\log p_0(\Lambda^c)$.
LLaDA's conditionals define no joint (sampled outputs depend on decoding order; TV $0.40$ in a two-token example), and the
per-position target $r_\varepsilon$ differs from the conditionals of the sequence-level projection: for the joint
$\{(f,a):.5,(b,f):.3,(b,c):.2\}$ the product of per-position targets puts $0.714$ on $(b,a)$, which has reference mass $0$.
Only the per-position target is used.
\end{remark}
\begin{remark}[Coverage constant]
For a training measure $\mu$ with $\mu>0$ on visited fact-free states, Proposition~\ref{prop:hazard}(d) gives
$\Prob_\pi(\Lambda_\F)\le C_{\pi|\mu}\E_\mu[H]$ with $C_{\pi|\mu}=\sup_s d^{\mathrm{ff}}_\pi(s)/\mu(s)$, where $H(s)=\sum_{i\in\M(s)}\min\{1,\rho(s,i)\}$
and $d^{\mathrm{ff}}_\pi(s)=\sum_n\Prob(s_n=s,\tau\ge n)$ is the pre-$\tau$ visitation \emph{mass} ($\le N$). For a mixture of normalised visitations,
$C\le|d^{\mathrm{ff}}_{\pi_k}|/(w_k(1-\alpha))$ ($\alpha$ = TF fraction of $\mu$), not $1/w_k$.
\end{remark}
\begin{remark}[Scale]
At LLaDA's default (\texttt{gen\_length} 128, block 32, 128 steps), $\sum_n|E_n|=2112$; certifying $5\%$ leakage at $T=1$ through
Lemma~\ref{lem:T}(iv) with $\varepsilon=0$ requires a mean per-position KL of order $2.4\times10^{-5}$. The certificate is therefore an
evaluation tool, not a consequence of training loss values.
\end{remark}
\begin{remark}[Floor $\varepsilon$]
The target's own fixed point leaks at least $1-(1-\varepsilon)^{n}$ under $T=1$ random remasking, $n$ = number of committed positions
with $p_0(\F)>\varepsilon$ ($0.8\%$ at $n=8$, $12\%$ at $n=128$ for $\varepsilon=10^{-3}$). With $\varepsilon=0$ the forward-KL
objective has no finite minimiser; in a 1-D Adam example the $\F$-logit offset keeps drifting sub-linearly ($-8.2,-15.8,-26.5$
after 5k/20k/80k steps), so the step budget is the strength knob.
\end{remark}
\begin{remark}[SFT weighting]
Under LLaDA's $1/t$-weighted SFT masking, the share of expected forget weight on (fact position, all fact positions masked)
pairs is $(1-\epsilon_t^{\,m})/((1-\epsilon_t)L)\approx 1/L$; this motivates dropping the $1/t$ weight but does not mean other
terms cannot change leakage (gradient ascent at a non-fact position can raise $p(\F)$ there).
\end{remark}
\end{document}
